\documentclass[10pt,a4paper]{article}
\usepackage{fontspec,amsmath,amssymb,graphicx}
\usepackage[margin=18mm]{geometry}
\setmainfont{Times New Roman}
\setlength{\parindent}{0pt}
\setlength{\parskip}{5pt}
\allowdisplaybreaks
\sloppy
\emergencystretch=2em
\begin{document}
\section*{VD\_4\_5\_2\_02 — Lời giải bổ sung}
\subsection*{Kiến thức và trình tự giải}
\textbf{Bước 1 — Phân biệt đạo hàm cấp hai, ba và bốn.} Quy ước $\widehat f(\xi)=\int e^{-i\xi x}f(x)dx$, biến đổi ngược có hệ số $1/(2\pi)$. Khi đó $\widehat{f^{(m)}}=(i\xi)^m\widehat f$, $\widehat{xf}=i\partial_\xi\widehat f$, $\widehat{xf_x}=-\widehat f-\xi\partial_\xi\widehat f$.

\textbf{Bước 2 — Giải phương trình của biến đổi.} $u_t=u_{xx}$ cho $\widehat u=e^{-t\xi^2}\widehat f$. $u_t=u_{xxx}$ cho $\widehat u=e^{-it\xi^3}\widehat f$. $u_t=-u_{xxxx}$ cho $\widehat u=e^{-t\xi^4}\widehat f$. Với $xu_x$, giải PDE cấp một theo $\xi$ bằng đường đặc trưng, gồm cả hệ số biên độ $-\widehat u$.

\textbf{Bước 3 — Biến đổi ngược và ý nghĩa tích phân.} Nhân bậc ba là $A_t(x)=(3t)^{-1/3}\operatorname{Ai}(-x/(3t)^{1/3})$, tích phân dao động. Nhân bậc bốn là $B_t(x)=(2\pi)^{-1}\int e^{i\xi x-t\xi^4}d\xi$, hội tụ tuyệt đối theo $\xi$ và có thể đổi dấu. Không áp dụng tính dương của Gauss hay nguyên lý cực đại cấp hai cho hai nhân này.

\textbf{Bước 4 — Miền hữu hạn cần kiểm tra biên và toán tử liên hợp.} Với $u_t=Ku_{xxx}$, $u(0)=u_x(0)=u(L)=0$, năng lượng cho $I'=-K u_x(L)^2/2$. Toán tử liên hợp là $-KD^3$ với $Y(0)=Y(L)=Y'(L)=0$. Đây không phải bài Sturm–Liouville tự liên hợp cấp hai; không suy mọi trị riêng đều thực và không tự dùng hệ số $\int fX/\int X^2$. Có thể giải bằng biến đổi Laplace theo thời gian và nghiệm cơ bản của ODE bậc ba, tránh khẳng định cơ sở chưa chứng minh.

\textbf{Bước 5 — Nêu lớp dữ liệu đủ cho công thức.} Với bậc ba trên toàn trục, công thức Fourier hợp lệ trước hết cho dữ liệu Schwartz hoặc $L^2$ theo nghĩa biến đổi Fourier. Dữ liệu liên tục bị chặn là phân phối tempered; phép nhân $e^{-it\xi^3}\widehat f$ vẫn hợp lệ vì các đạo hàm của thừa số tăng đa thức, nên xác định nghiệm trong $\mathcal S\prime$. Tuy nhiên điều đó không bảo đảm tích phân Airy hội tụ tuyệt đối hay nghiệm tự động là hàm cổ điển. Ghi đúng ý nghĩa của công thức theo lớp dữ liệu.
\subsection*{Lời giải áp dụng cho bài}
\textbf{Bước 1. Năng lượng bậc ba và tính duy nhất.}
Với $K>0$, $u_t=Ku_{xxx}$, $u(0)=u_x(0)=u(L)=0$, hai lần tích phân từng phần cho
\[I'=K\int uu_{xxx}
=K[uu_{xx}]_0^L-\frac K2[u_x^2]_0^L
=-\frac K2u_x(L,t)^2\le0.\]
Bản in bỏ hệ số $1/2$ ở đẳng thức cuối; dấu suy giảm vẫn đúng. Áp dụng cho hiệu hai nghiệm có mặt đầu $0$, suy $I=0$, tức duy nhất trong lớp nghiệm đủ trơn/năng lượng thích hợp.

\textbf{Bước 2. Vì sao không dùng cơ sở trực giao cấp hai.}
Tách biến $u=XT$ cho $X'''=\lambda X$, $T'=K\lambda T$ và ba biên $X(0)=X'(0)=X(L)=0$. Với hàm riêng phức, lấy phần thực của tích phân với $\overline X$:
\[\operatorname{Re}\lambda\int_0^L|X|^2dx=-\frac12|X'(L)|^2\le0.\]
Điều này không chứng minh $\lambda$ thực. Toán tử $D^3$ không tự liên hợp; liên hợp là $-D^3$ với biên $Y(0)=Y(L)=Y'(L)=0$, suy từ hạng $[X''\overline Y-X'\overline{Y'}+X\overline{Y''}]_0^L$. Vì vậy chưa được khẳng định dãy hàm riêng thực trong sách là cơ sở đầy đủ.
Nếu chỉ xét nhánh thực $\lambda=-k^3$, nghiệm thỏa hai biên trái là
\[X=e^{-kx}-e^{kx/2}\cos(\sqrt3kx/2)+\sqrt3e^{kx/2}\sin(\sqrt3kx/2)
=e^{-kx}+2e^{kx/2}\sin(\sqrt3kx/2-\pi/6).\]
Biên phải cho $e^{-3kL/2}=-2\sin(\sqrt3kL/2-\pi/6)$. Dấu trước sine trong $X$ phải là $+$; dấu $-$ ở công thức hàm riêng của bản in làm $X(0)=2$, sai biên.

\textbf{Bước 3. Giải đầy đủ bằng Laplace, không giả định cơ sở chưa chứng minh.}
Đặt $U(x,p)=\int_0^\infty e^{-pt}u(x,t)dt$, $\operatorname{Re}p>0$. Từ $pU-f=KU'''$:
\[U'''-\frac pK U=-\frac fK,\quad U(0,p)=U_x(0,p)=U(L,p)=0.\]
Đặt nghiệm cơ bản
\[S_p(x)=\sum_{m=0}^\infty\frac{(p/K)^m x^{3m+2}}{(3m+2)!}.\]
Đạo hàm chuỗi cho $S_p'''=(p/K)S_p$, $S_p(0)=S_p'(0)=0,S_p''(0)=1$. Biến thiên hằng số cho
\[U(x,p)=C(p)S_p(x)-\frac1K\int_0^xS_p(x-y)f(y)dy.\]
Đạo hàm tích phân ba lần tạo đúng $f(x)$ nhờ $S_p''(0)=1$. Biên $U(L)=0$ xác định
\[C(p)=\frac{\int_0^LS_p(L-y)f(y)dy}{KS_p(L)}.\]
Mẫu không bằng $0$ khi $\operatorname{Re}p>0$: nếu $S_p(L)=0$, $S_p$ là hàm riêng không tầm thường với trị riêng $p/K$ có phần thực dương, trái đẳng thức năng lượng.

\textbf{Bước 4. Công thức nghiệm và kiểm tra.}
\[
u(x,t)=\mathcal L^{-1}_{p\to t}
\left\{\frac{S_p(x)}{KS_p(L)}\int_0^LS_p(L-y)f(y)dy
-\frac1K\int_0^xS_p(x-y)f(y)dy\right\}.
\]
Ở đây $\mathcal L^{-1}H=(2\pi i)^{-1}\int_{\gamma-i\infty}^{\gamma+i\infty}e^{pt}H(p)dp$, $\gamma>0$, hiểu theo phép đảo Laplace trong lớp năng lượng; với dữ liệu trơn tương thích có nghiệm cổ điển. Công thức $U$ thỏa ODE, cả ba biên và hạng mặt đầu $f$ trong $pU-f=KU'''$, nên phép đảo trả lại đúng bài toán. Với $f\in L^2$, nghiệm được hiểu là nửa nhóm co sinh bởi $KD^3$ trên miền $H^3$ có ba biên trên; tiêu tán vừa chứng minh và nghiệm resolvent cho mọi $p>0$ bảo đảm tồn tại và duy nhất. Đây là công thức giải cho dữ liệu tổng quát mà không bỏ các thành phần phổ phức hoặc giả định các hệ số trực giao.\end{document}